\documentclass[a4paper, 11pt]{article} % das Papierformat zuerst
\usepackage[utf8]{inputenc}
\usepackage{geometry}
\geometry{a4paper, top=25mm, left=35mm, right=25mm, bottom=25mm}
\usepackage{graphicx}
\usepackage{color}
\usepackage{epstopdf}
\usepackage[T1]{fontenc}
\usepackage{setspace}
\usepackage{tabularx}
\usepackage{blindtext}
\usepackage[ngerman]{babel} %deutsche Silbentrennung
\usepackage{titlesec}
\usepackage{enumitem}
\usepackage{ifthen}
\usepackage[colorlinks=true,linkcolor=fsi]{hyperref}

\definecolor{fsi}{RGB}{0,81,150}

\newcounter{summe}

\newcommand{\abstimmung}[4]{
	\setcounter{summe}{#3}
	\addtocounter{summe}{#4}
	\begin{flushleft}
		#1\\
	Es wird über den Antrag abgestimmt.
	\end{flushleft}
	\ifthenelse{\equal{#3}{0}\AND\equal{#4}{0}}{
	\begin{center}
		\textbf{Der Antrag wird einstimmig angenommen.}
	\end{center}
	}{
	\begin{center}
		#2 \ifthenelse{\equal{#2}{1}}{Stimme}{Stimmen} dafür, #3 \ifthenelse{\equal{#3}{1}}{Stimme}{Stimmen} dagegen, #4 \ifthenelse{\equal{#42}{1}}{Enthaltung}{Enthaltungen}\\
		\ifthenelse{#2>\value{summe}}{
		\textbf{Der Antrag ist somit angenommen.}
		}{
		\textbf{Der Antrag ist somit abgelehnt.}
		}
	\end{center}
	}
}


% VARIABLEN
\newcommand{\protokoller}{Johannes Kirchner}
\newcommand{\dateOfMeeting}{22. Juni 2015}
\newcommand{\TeXer}{Johannes Kirchner}
\newcommand{\fsiPresident}{Georg Schäfer}
\newcommand{\protocolID}{8}

\begin{document}
%deckblatt start

\doublespacing
\thispagestyle{empty}

\begin{center}
\includegraphics[width=10.0cm]{../logo_fsi.eps}

\vspace*{\fill}
{\LARGE \textbf{Protokoll der Sitzung des Fachschaftsrates \\vom \dateOfMeeting}}\\
Fachschaftsrat Informatik\\
Trier University of Applied Sciences\\
\vspace{2.5cm}
\textit{
	Protokollführer: \textbf{\protokoller} \\
	\LaTeX - Umsetzung von \TeXer\\
	am \today\\
}
Protokoll ID: \textbf{\#\protocolID}
\vfill
\end{center}

\hspace*{-35cm}
\textcolor{fsi}{\rule{64.9cm}{15pt}}
\pagebreak

\setcounter{tocdepth}{2}
\tableofcontents
\pagebreak

\section{Eröffnung}
Als Protokollführer wird \protokoller~bestimmt.\\
Der Sprecher des Fachschaftsrates \fsiPresident~eröffnet die Sitzung um 19:23 Uhr.
\\\\
\textbf{Es wird festgestellt, dass der Fachschaftsrat beschlussfähig ist.}\\
\textbf{Es fehlen: Santo Pfingsten}
\\\\
Das Protokoll der letzten Sitzung wurde verlesen und genehmigt.
%Die Genehmigung des Protokolls der letzten Sitzung wird auf die nächste Sitzung vertagt.

%\abstimmung{Es wird der Antrag gestellt, Dinge zu tun}{5}{6}{2} -- {dafür}{dagegen}{enthaltungen}

\section{Events}
Der Termin für die GDW ist bereits fesgelegt und es wurden Räume reserviert. 3 Leute machen die GDW als Fachseminar.
Tobias Meier und Santo Pfingsten werden sich mit Janine Kleinbauer treffen, um das Abschlussgrillen des 100-Tage-Projektes zu planen.
Fabio plant, die nächste NOTLD noch vor der Reading-Week zu veranstalten.
Der nächste Spieleabend ist am 09.07, im Fachschaftsraum befinden sich von der letzten noch Spiele, den Besitzer unbekannt ist.

\section{NIV}
Das Geld für den Server wurde endlich überwiesen.
Das Banner kommt morgen. Es wir überlegt Rabatt zu beantragen, da erst verspätet geliefert wurde.
Nachdem die Anträge für Mehrfachstecker genehmigt wurden, wurden diese bereits bestellt.
Es sind 34 Kästen Mate bei Herresthal bestellt.
Wie die Medienanlage der Aula bedient wird, wird beim Aufbau geklärt.
Der Fachschaftsrat der Uni konnte die Plakate nicht aufhängen.
Bis jetzt gibt es 51 Reservierungen.
Für Auf-/Abbau und verkauf gibt es bisher zu wenige Zusagen im Doodle.

\section{Finanzen}
Wir heben ein neues Kassenbuch und Datumsstempel. In der Ablage befindet sich eine Liste für Anschaffungen, bald soll auch auf der Webseite ein Formular dazu eingerichtet werden.
Dringend benötigt werden (AA-)Batterien.
In den letzten Tagen ist ein leerer Kasten Mate verschwunden.
Zu Schrank und sofa gibt es noch nix neues, Dominik schickt ein Angebot für einen neuen Schrank an den Rat.
Alle Anträge, die noch bei AStA lagen, sind genehmigt worden.
Marco schlaegt vor, einen neuen Beamer zu kaufen. Es wird vorgeschlagen, Oechsle / den Fachbereich zu bitten, das alte Leihgerät auszutauschen.

\section{Raum}
Nicht-rats-mitglieder haben noch keinen Raumzugang.
Der Putzplan klappt nicht, wer seinen Dienst nicht tut, ist schwer zu kontrollieren.
Die alten Lacktische wurden durch neue ersetzt.
Die Pflanze wird unregelmäßig und zu selten gegossen, lebt aber noch.

\section{Sonstiges}
\abstimmung{Es wird der Antrag gestellt, ab jetzt die Beschluesse durchzunummerieren}{12}{0}{0}

Trotz der Unfähigkeit des FSR Info der Uni Trier, bei sich Werbung für uns zu machen, wird weiterhin Werbung für diesen an der Hochschule verteilt. Georg setzt sich stark dafuer ein. Stefan ebenso.

\abstimmung{Es wird der Antrag gestellt, keine Werbung mehr fuer den Fachschaftsrat Informatik der Uni an der Hochschule zu machen}{2}{6}{4}

\fsiPresident~schliesst sitzung um 19:57

\pagebreak
\section{Bestätigung des Protokolls}
Der Sprecher des Fachschaftsrates Informatik sowie der Protokollführer dieses Protokolls bestätigen mit Ihrer Unterschrift unter diesem Protokoll, dass selbiges inhaltlich korrekt ist und alle darin aufgeführten Beschlüsse so wie beschrieben vom Rat getragen werden.
\\

\vspace{3.5cm}
\hrulefill \hfill \hrulefill

\fsiPresident \hfill \protokoller

{\footnotesize (Sprecher)\hfill (Protokollführer)}
\end{document}
